\documentclass{article}

\usepackage[table,xcdraw]{xcolor}
\usepackage{amsmath,amssymb,fullpage,adjustbox}
\usepackage{dsfont}
\setlength\parindent{0pt}
\begin{document}

	\section*{Ancestral skies Grounds keeping}
	By Harry Hawkins 48019675
	
	\section*{General formulation}
	The variables, definitions and transition function used for all three models.
	The specififc values of $\mathds{P}(s | C_{st})$ will be adressed for each model along with the value function used.

	\subsection*{Sets}
	\begin{tabular}{rl}
		$S$ & Sites\\
		$P$ & Species
	\end{tabular}

	\subsection*{Data}
	\begin{tabular}{rl}
		$L_s$ & The species found at site $s\in S$\\
		$N_s$ & The neighbouring sites of site $s\in S$\\
		$\mathds{P}(s | C_{st})$ & Probability of site $s$ perishing by $t+1$ given the state $C_{st}$\\
		$\mathds{P}(s' | s, A_s)$ & The transition function given as GenerateOutcomes giving all possible states\\
		& when all possible actions are applied and the probability of them occuring.
		
	\end{tabular}
		

	\subsection*{Stages}
	\begin{tabular}{rl}
		$T$ & Months
	\end{tabular}

	\subsection*{States}
	\begin{tabular}{rl}
		$C_{st}$ & Condition of site $s\in S$ at the end of month $t\in T$\\
		$x_t$ & The collection of $C_{st}$ $\forall s\in S$
	\end{tabular}

	\subsection*{Actions}
	\begin{tabular}{rl}
		$A_s$ & Save site $s\in S$, limited to one per month, by setting $C_{st} = 1$
	\end{tabular}

	\section*{Models}
	For notation purposes, define $SpeciesSaved$ as folows:
	\begin{equation*}
		\sum_{p\in P} \begin{cases}
			1 & \text{if $\exists s\in S$ $|$ $C_{st} = 1$ and $p\in L_s$}\\
			0 & \text{otherwise}
		\end{cases}
	\end{equation*}

	\subsection*{Com13}
	Begin with initial state of all site are fragile.\\
	$\mathds{P}(s | C_{st}) = 0.2$ $\forall s\in S$\\
	\begin{equation*}
	 	V(x_t) = \begin{cases}
			SpeciesSaved & \text{if no fragile sites remain}\\
			\max{} \sum \mathds{P}(x_{t+1})V(x_{t+1}), x_{t+1} = x_t + A_s & \text{otherwise}\\
		\end{cases}
	\end{equation*}

	\subsection*{Com14}
	Begin with initial state of all site are fragile.\\
	The new probability of a site perishing is dependent on the neighbouring sites state.\\
	$\mathds{P}(s | C_{st}) = 0.2 + 0.05 * \sum_{N_s} C_{st} = -1$ $\forall s\in S$\\
	\begin{equation*}
	 	V(C_{st}) = \begin{cases}
			SpeciesSaved & \text{if no fragile sites remain}\\
			\max{} \sum \mathds{P}(x_{t+1})V(x_{t+1}), x_{t+1} = x_t + A_s & \text{otherwise}\\
		\end{cases}
	\end{equation*}

	\subsection*{Com15}
	$\mathds{P}(s | C_{st}) = 0.2 + 0.05 * \sum_{N_s} C_{st} = -1$ $\forall s\in S$\\
	\begin{equation*}
		V(x_t) = \begin{cases}
			1 & \text{no fragile states remain and all required species are saved}\\
			0 & \text{no fragile states remain and not all required species are saved}\\
			\max{} \sum \mathds{P}(x_{t+1})V(x_{t+1}), x_{t+1} = x_t + A_s & \text{otherwise}
		\end{cases}
	\end{equation*}
	A note on the final commmunication: The first attempt of the model produced a schedual that did not 
	include all sites. To adjust it the probability was skewed to favour higher sites. The probability value did
	not change the result up to 7 significant values.

	\newpage
	\subsection{AI Decleration}
	AI was used in communication 15 to adjust the model to favour full scheduals. See below
	\\Q. This code skips one of 9 sites in the solution to find the optimal probability of 
	protecting the species given in survives(s). I need it to favour producing a schedual with all 9 sites.
	How do you suggest i do this. Do not show me the code solution, just need the idea.
	$_Pbs = {}\\
	def com15(s):\\
    if s in _Pbs:\\
        return _Pbs[s]\\
    fragile = [i for i in Sites if s[i] == 0]\\
    if fragile == []:\\
        if survives(s):\\
            _Pbs[s] = (1, [])\\
            return (1, [])\\
        _Pbs[s] = (0, [])\\
        return(0, [])\\
    best_sched = []\\
    best_prob = 0\\
    for a in fragile:\\
        copy_s = list(s)\\
        copy_s[a] = 1\\
        copy_s = tuple(copy_s)\\
        outcomes = GenerateOutcomes(copy_s, probabilities(copy_s))\\
        temp = 0\\
        sched = []\\
        best_temp = -1\\
        for y in outcomes:\\
            prob, sch = com15(y[1])\\
            temp += y[0]*prob\\
            if prob > best_temp:\\
                best_temp = prob\\
                sched = list(sch)\\
        if temp > best_prob:\\
            best_prob = temp\\
            best_sched = [a] + sched\\
    _Pbs[s] = (best_prob, best_sched)\\
    return (best_prob, best_sched)\\
	$
\\
	A. It produced 3 cases but I only applied the second.\\
	"Modify the objective to:
		$score = probability + \epsilon \cdot  schedule_length$

	This is a standard trick in DP when you want to bias toward longer or 
	shorter solutions without changing the primary objective."

\end{document}